\textbf{Phần 1. Trắc nghiệm lựa chọn}
\setcounterpageref{ex}{0}

\begin{ex} %[2D14N68]
Số gồm $5$ trăm, $8$ chục và $3$ đơn vị viết là
\choice
{\True $583$}
{$538$}
{$853$}
{$385$}
\loigiai{
Số đó là $583$.\\
\textbf{Đáp án của câu hỏi là A}
}
\end{ex}

\begin{ex} %[2D14N68]
Số $905$ đọc là
\choice
{Chín trăm năm mươi}
{\True Chín trăm linh năm}
{Chín trăm năm}
{Chín không năm}
\loigiai{
Đọc là Chín trăm linh năm.\\
\textbf{Đáp án của câu hỏi là B}
}
\end{ex}

\begin{ex} %[2D14H68]
Số liền trước của số $800$ là
\choice
{$801$}
{$790$}
{\True $799$}
{$700$}
\loigiai{
$800 - 1 = 799$.\\
\textbf{Đáp án của câu hỏi là C}
}
\end{ex}

\begin{ex} %[2D14H68]
Số liền sau của số $999$ là
\choice
{$998$}
{\True $1\,000$}
{$1001$}
{$990$}
\loigiai{
$999 + 1 = 1000$.\\
\textbf{Đáp án của câu hỏi là B}
}
\end{ex}

\begin{ex} %[2D14H68]
Số $672$ được viết thành tổng các trăm, chục, đơn vị là
\choice
{$60 + 70 + 2$}
{$600 + 7 + 2$}
{\True $600 + 70 + 2$}
{$600 + 72$}
\loigiai{
$672 = 600 + 70 + 2$.\\
\textbf{Đáp án của câu hỏi là C}
}
\end{ex}

\begin{ex} %[2D14H68]
Điền dấu thích hợp: $856 \dots 865$.
\choice
{\True $<$}
{$>$}
{$=$}
{Không so sánh được}
\loigiai{
$856 < 865$.\\
\textbf{Đáp án của câu hỏi là A}
}
\end{ex}

\begin{ex} %[2D14V68]
Trong các số $645, 564, 456, 654$, số lớn nhất là
\choice
{$645$}
{\True $654$}
{$564$}
{$456$}
\loigiai{
Số lớn nhất là 654.\\
\textbf{Đáp án của câu hỏi là B}
}
\end{ex}

\begin{ex} %[2D14V68]
Trong các số $789, 879, 798, 897$, số bé nhất là
\choice
{$798$}
{$879$}
{$897$}
{\True $789$}
\loigiai{
Số bé nhất là 789.\\
\textbf{Đáp án của câu hỏi là D}
}
\end{ex}

\begin{ex} %[2D14V68]
Sắp xếp các số $312, 123, 231, 132$ theo thứ tự từ bé đến lớn:
\choice
{\True $123, 132, 231, 312$}
{$312, 231, 132, 123$}
{$123, 231, 132, 312$}
{$132, 123, 231, 312$}
\loigiai{
$123 < 132 < 231 < 312$.\\
\textbf{Đáp án của câu hỏi là A}
}
\end{ex}

\begin{ex} %[2D14V68]
Số chẵn lớn nhất có 3 chữ số khác nhau là
\choice
{$998$}
{$988$}
{$978$}
{\True $986$}
\loigiai{
Số đó là 986.\\
\textbf{Đáp án của câu hỏi là D}
}
\end{ex}

\begin{ex} %[2D14V68]
Số tròn chục bé nhất có ba chữ số khác nhau là
\choice
{$100$}
{$102$}
{\True $120$}
{$110$}
\loigiai{
Số đó là 120.\\
\textbf{Đáp án của câu hỏi là C}
}
\end{ex}

\begin{ex} %[2D14C68]
Tìm số có 3 chữ số biết chữ số hàng trăm là số lớn nhất có một chữ số, chữ số hàng chục bằng một phần ba chữ số hàng trăm và chữ số hàng đơn vị bằng 0.
\choice
{$960$}
{$910$}
{$900$}
{\True $930$}
\loigiai{
Hàng trăm: 9. Hàng chục: $9 : 3 = 3$. Hàng đơn vị: 0. Số đó là $930$.\\
\textbf{Đáp án của câu hỏi là D}
}
\end{ex}

\textbf{Phần 2. Trắc nghiệm đúng sai}
\setcounterpageref{ex}{0}

\begin{ex} %[2D14H68]
Xét tính đúng sai về số trong phạm vi 1 000:
\choiceTF
{\True $1\,000$ là số bé nhất có bốn chữ số.}
{\True Số liền trước của số 600 là 599.}
{Số 408 viết thành tổng là $400 + 80$.}
{\True Số 408 viết thành tổng là $400 + 8$.}
\loigiai{
\begin{itemize}
    \item 1000 bé nhất có 4 chữ số: Đúng.
    \item Liền trước 600 là 599: Đúng.
    \item $408 = 400 + 8 \neq 400 + 80$. Ý 3 Sai, Ý 4 Đúng.
\end{itemize}
\begin{itemize}
    \item a) \textbf{Đúng}.
    \item b) \textbf{Đúng}.
    \item c) \textbf{Sai}.
    \item d) \textbf{Đúng}.
\end{itemize}
}
\end{ex}

\begin{ex} %[2D14V68]
Cho các số: $450, 540, 405, 504$. Xét tính đúng sai:
\choiceTF
{\True Số lớn nhất là 540.}
{\True Số bé nhất là 405.}
{Sắp xếp theo thứ tự giảm dần là $540 > 450 > 504 > 405$.}
{\True Sắp xếp giảm dần đúng là $540 > 504 > 450 > 405$.}
\loigiai{
\begin{itemize}
    \item Lớn nhất 540: Đúng.
    \item Bé nhất 405: Đúng.
    \item $504 > 450$. Ý 3 Sai, Ý 4 Đúng.
\end{itemize}
\begin{itemize}
    \item a) \textbf{Đúng}.
    \item b) \textbf{Đúng}.
    \item c) \textbf{Sai}.
    \item d) \textbf{Đúng}.
\end{itemize}
}
\end{ex}

\textbf{Phần 3. Trả lời ngắn}
\setcounterpageref{ex}{0}

\begin{ex} %[2D14H68]
Viết số gồm 7 trăm, 0 chục và 9 đơn vị.
\shortans{709}
\loigiai{
\textbf{Đáp số:} $709$\\
709.
}
\end{ex}

\begin{ex} %[2D14H68]
Tìm số liền trước của số $900$.
\shortans{899}
\loigiai{
\textbf{Đáp số:} $899$\\
$900 - 1 = 899$.
}
\end{ex}

\begin{ex} %[2D14V68]
Tìm số liền sau của số $889$.
\shortans{890}
\loigiai{
\textbf{Đáp số:} $890$\\
$889 + 1 = 890$.
}
\end{ex}

\begin{ex} %[2D14V68]
Viết kết quả của tổng: $800 + 60 + 4$.
\shortans{864}
\loigiai{
\textbf{Đáp số:} $864$\\
864.
}
\end{ex}

\begin{ex} %[2D14V68]
Tìm số lớn nhất có 3 chữ số khác nhau.
\shortans{987}
\loigiai{
\textbf{Đáp số:} $987$\\
987.
}
\end{ex}

\begin{ex} %[2D14C68]
Từ 3 chữ số $0, 4, 7$, viết số bé nhất có 3 chữ số khác nhau.
\shortans{407}
\loigiai{
\textbf{Đáp số:} $407$\\
407.
}
\end{ex}
